\begin{example}
	Существует ли множество, неизмеримое по Лебегу? К сожалению, да. Рассмотрим отрезок $[0; 1]$ и введём отношение эквивалентности на нём:
	\[
		x \sim y \lra (x - y) \in \Q
	\]
	Так как это действительно отношение эквивалентности, то отрезок $[0; 1]$ разбивается на классы эквивалентности:
	\[
		[0; 1] = \bscup_{\alpha} K_\alpha
	\]
	Пусть $E$ - это множество, которое образовано из элементов различных классов (по одному из каждого. Здесь мы воспользовались аксиомой выбора). Занумеруем все рациональные числа отрезка $[-1; 1]$:
	\[
		[-1; 1] \cap \Q = \{r_n\}_{n = 1}^\infty
	\]
	Предположим, что $E$ измеримо по Лебегу. Тогда
	\[
		\forall m \in \N\ \ E_m := \{y \colon y = x + r_m, x \in E\}
	\]
	тоже измеримы по лебегу, причём $\forall m \in \N\ \mu(E_m) = \mu(E)$.
	
	Докажем, что $\forall k \neq j\ E_k \cap E_j = \emptyset$. От противного, пусть $y_0 \in E_k \cap E_j$. В таком случае
	\[
		y_0 = x_k + r_k = x_j + r_j,\ x_k \in E,\ x_j \in E
	\]
	Если переместить иксы в одну сторону, то получится интересный факт:
	\[
		x_j - x_k = r_k - r_j \in \Q \lra x_j \sim x_k
	\]
	Но такого быть не может по определению множества $E$. Теперь, возьмём произвольный $y \in [0; 1]$. Тогда $\exists \alpha_0 \such y \in K_{\alpha_0}$. Более того, $\exists x \in E \cap K_{\alpha_0}$. Стало быть
	\[
		\exists m \in \N \such y = x + r_m \Ra y \in E_m
	\]
	В итоге, мы имеем следующую цепочку вложений и можем воспользоваться свойствами мер:
	\[
		[0; 1] \subset \bscup_{m = 1}^\infty E_m \subset [-1; 2] \Lora 1 \le \sum_{m = 1}^\infty \mu(E_m) = \sum_{m = 1}^\infty \mu(E) \le 3
	\]
	Бесконечная сумма одинаковых неотрицательных чисел должна быть зажата среди двух положительных, а такого быть не может.
\end{example}

\begin{note}
	Вообще говоря весь анализ ранее (за исключением примера всюду плотной функции) существенно опирался только на так называемую счётную аксиому выбора, иначе говоря, её нельзя считать неверной, если мы хотим иметь адекватную науку. Но при этом её можно усилять, в частности есть два достаточно популярных взаимоисключающих варианта: аксиома выбора и аксиома детерменированности. Первая старше и именно она используется большинством, приводя, в том числе к существованию неизмеримого по Лебегу множества. Вторая же наоборот, позволяет доказать, что все множества по Лебегу измеримы, но касаться этой части науки мы здесь не будем.
\end{note}

\begin{note}
	В примере выше вместо $[0; 1]$ можно взять любое множество $A \subset \R \colon 0 < \mu(A) < \infty$. Просто в паре мест понадобится заменить $[0; 1]$ на $A$. Другими словами, это означает, что из любого множества конечной меры можно выбрать неизмеримое подмножество.
\end{note}

\subsection{Отображения измеримых множеств}

\begin{definition}
	Отображение $\vv{f}: \D\to \R^m,\ \D\subset\R^n$, называется \textit{(непрерывно) дифференцируемым во внутренней точке} $\vv{x}\in \D$, если $\vv{f} = (f_1,...,f_m),\ f_j,\ j=1,...,m$ (непрерывно) дифференцируемы в точке $\vv{x}$. Производной $\vv{f}(\vv{x})$ назовём матрицу Якоби:
	\[
		\vv{f}':= \left(\pd{f_j}{x_k}(\vv{x}) \right)^{m,\  n}_{j=1,\ k=1} = \begin{pmatrix}
			\grad{f_1(\vv{x})}\\
			\vdots\\
			\grad{f_m(\vv{x})}
		\end{pmatrix}
	\]
	Взаимно однозначное дифференцируемое неособое (т.е. $\det{\vv{f}\neq 0}$) отображение области $\D$ в $\R^n$ называется \textit{диффеоморфизмом}.
\end{definition}


\begin{example}
	$m=1$, тогда: 
	\[
	f'(\vv{x}) = \grad{f(\vv{x})}
	\]
\end{example}

\begin{example}
	$n=1$, тогда: 
	\[
		\vv{f}'(x) = \begin{pmatrix}
			f_1'(x)\\
			\vdots\\
			f_m'(x)
		\end{pmatrix}
	\]
\end{example}

\begin{example}
	\textit{(Правило дифференцирования сложной функции)}:
	\[
		h = f \circ \vv{g}
	\]
	Тогда:
	\[
		h'= f'(\vv{g})\cdot \vv{g}'
	\]
\end{example}

\begin{notation}
	$|\vv{f}'(\vv{x})|:=\max\limits_{1\le j\le m}{|\grad{f_j(\vv{x})}|}$
\end{notation}

\begin{note}
	Вообще говоря, последнее обозначение является костылём, для того чтобы не вводить теорию нормы и нормированных пространств. И на самом деле собственно должно обозначать матричную норму.
\end{note}

\begin{example}
	А теперь пусть композиция двух вектор-функций с векторной областью определения:
	\[
		\vv{h}=\vv{f}\circ\vv{g}
	\]
	На самом деле ничего не поменяется:
	\[
	\vv{h}'= \vv{f}'(\vv{g})\cdot \vv{g}'
	\]
	(Проверяется это видимо внимательным взглядом)
\end{example}

\begin{note}\textit{(Общие сведения)}
	\begin{itemize}
		\item $E$ - измеримо(по Лебегу), тогда $\mu(E+\vv{x})=\mu(E)$ (т.к. сдвиги не меняют объём элементарных множеств)
		\item $\alpha E:={\alpha \vv{x} : \vv{x}\in E}$, тогда из очевидности этого свойства для элементарных получим ещё и $\mu(\alpha E) = \alpha^n\mu(E)$
	\end{itemize}
\end{note}

\begin{note}
	Последнее, являющееся растяжением по каждой оси в $\alpha$ раз, можно обобщить в расширение в разное количество раз по разным осям, что выразится, как домножение на соответствующую матрицу, а норма, в свою очередь, будет домножаться на определитель этой матрицы. Позже мы узнаем, что аналогично домножать можно вообще на любые матрицы.
\end{note}

\begin{definition}
	Множество $\D\subset\R^n$ называется выпуклым, если $(\forall \vv{x}, \vv{y}\in \D)\emb[\vv{x}, \vv{y}]\subset\D$, т.е. $(\forall t\in[0, 1])\emb t\vv{x} + (1-t)\vv{y}\in\D$. 
\end{definition}

\begin{lemma}
	(Теорема Лагранжа для отображений) Если $\vv{f}:\Omega\to\R^m$ непрерывно дифференцируемо на ограниченном открытом множестве $\Omega\supset\cl{\D}$, $\D\subset\R^n$ выпукло, то $(\forall \vv{x}, \vv{y}\in\D)\emb |\vv{f}(\vv{x}) - \vv{f}(\vv{y})|\leq \max\limits_{\vv{\xi}\in\cl{\D}}{|\vv{f}'(\vv{\xi})|\sqrt{m}|\vv{x}-\vv{y}|}$ Здесь $\xi$ внезапно от $\vv{x}$ и $\vv{y}$ не зависит.
\end{lemma}

\begin{proof}
	Так как координатные фунуции $f_i$ дифференцируемы на $\Omega$ то по формуле Тейлора с остаточным членом в форме Лагранжа
	\[
		(\forall \vv{x}, \vv{y}\in\D)\ (\exists\vv{\xi}\in[\vv{x}, \vv{y}]) \emb f_i(\vv{x}) - f_i(\vv{y}) = df_i(\vv{\xi})
	\]
	Отсюда
	\[
		|f_i(\vv{x}) - f_i(\vv{y})| = |\trbr{\grad{f_i(\vv{\xi})}, \vv{x} - \vv{y}}| \leq |\grad{f_i(\vv{\xi})}||\vv{x} - \vv{y}|
	\]
	Значит какие бы мы не взяли $\vv{x}$ и $\vv{y}$ модуль разности между значениями в них можно оценить с помощью наибольшего значения функции $|\grad{f_i(\vv{\xi})}|$ на всём множестве, но чтобы максимум было брать корректно мы добьём $\D$ до компакта, взяв у него замыкание, как у ограниченного множества.
	\[
		|\vv{f}(\vv{x}) - \vv{f}(\vv{y})| = \sqrt{\sum\limits_{i = 1}^m |f_i(\vv{x}) - f_i(\vv{y})|^2} \leq \max\limits_{\vv{\xi}\in\cl{\D}}{|\vv{f}'(\vv{\xi})|\sqrt{m}|\vv{x}-\vv{y}|}
	\]
\end{proof}

\begin{theorem}(Диффеоморфный образ куба меньшей размерности)
	Если $\vf: \Omega\to\R^m$ - непрерывно дифференцируема на открытом множестве $\Omega\supset K_I$ - единичный куб в $\R^n$, $n<m$, то
	\[
		\jm(\vf(K_I))=0
	\]
\end{theorem}

\begin{proof}
	Разобьём $K_I$ на $N^n$ кубов со стороной длины $\frac{1}{N} = \delta$. По прошлой лемме, если пронумеровать кубы:
	\[
		(l = 1,...,N^n)\ (\forall\vx, \vy\in K_{\delta, l})\emb|\vf(\vx)-\vf(\vy)| \leq M\sqrt{m}|\vx - \vy|, \text{где}
	\]
	$M=\max\limits_{\vv{t}\in K_I}{|\vf'(t)|}$ \\
	Отсюда оценим меру образа одного кубика, просто упихнув его в кубик с ребром в два раза большим, чем полученная верхняя оценка на расстояние между точками.
	\[
		\upjm(\vf(K_{\delta, l}))\leq (2M\sqrt{mn}\delta)^m
	\]
	Значит
	\[
	\upjm(\vf(K_{I}))\leq N^n \upjm(\vf(K_{\delta, l}))\leq (2M\sqrt{mn})^m N^{n-m}
	\]
	Последнее, в силу отрицательности $n-m$, при стремлении $N$ к бесконечности, стремится к нулю. Что и требовалось доказать ибо из нулевой внешней меры, следует и обычная.
\end{proof}

\begin{theorem}
	(Диффеоморфный образ измеримого по Лебегу множества) Если $\vf: \Omega\to\R^n$ - диффеоморфизм на ограниченном открытом множестве $\Omega\supset\cl{\D}$, $\D\subset\R^n$-измеримо по Лебегу, то $\vf(\D)$ измеримо по Лебегу.
\end{theorem}

\begin{proof}
	Прежде всего докажем измеримость $\vf(P)$, где $P\subset\Omega$ - брус. $\vf(P) = \vf(\cl{P}) \backslash\vf(\vdelta P\backslash P)$ Аргумент уменьшаемого в правой части - компакт, его непрерывный образ тогда тоже, а значит слагаемое измеримо. Аргумент вычитаемого есть часть границы бруса, то есть подмножество граней бруса, которые можно покрыть кубами меньшей размерности. Это по прошлой теореме значит, что вычитаемое имеет меру нуль. Значит левая часть - разность измеримого и множества с мерой нуль, т.е. измеримо. \\
	\\
	Теперь воспользуемся критерием измеримости. По нему $(\forall\eps>0)\ (\exists N_\eps -\text{элементарное})\emb\mu(N_\eps\tr\D) < \eps$. В силу компактности $\D$ его можно покрыть конечным числом окрестностей-кубов для точек из $\Omega$. Пересекая с этим покрытием $N_\eps$ мы не уменьшим симметрическую разность с $\D$. Поэтому БОО $N_\eps\subset\Omega$.\\
	\\
	По доказательству прошлой теоремы $\mu^*(\vf(P))\leq(2Mn)^n|P|,\ M=\max\limits_{t\in\cl{\D}\cup\cl{N_{\eps_0}}}{|\vf(t)|}$. Здесь $N_{\eps_0}$ строится как элементарное множество в которое попадёт любое наше, из св-в компакта, как в прошлом абзаце.\\
	Итак, полученное неравенство верно для любого бруса, значит и для любого элементарного множества, значит и для любого измеримого множества:
	\[
		\mu^*(\vf(A))\leq(2Mn)^n\mu^*(A)
	\]
	В силу того, что мы внутри компакта, непрерывные образы открытых открыты, тогда по теореме о структуре открытых:
	\[
		\vf(\Int {P})=\bigsqcup\limits_{k=1}^\infty Q_k,\ \sum\limits_{k=1}^\infty|Q_k| <\infty
	\]
	Тогда заметим, что
	\[
		(\forall\delta>0)\ (\exists K)\emb \sum\limits_{k=K+1}^\infty|Q_k|\ <\delta \Rightarrow (\forall P\subset\Omega)\ (\exists M -\text{эл.}\subset\vf(P))\emb\mu(\vf(P)) - |M| < \delta
	\]
	Значит $(\forall \delta > 0)\ (\exists M_{\eps}\subset\vf(N_\eps))\emb \mu(\vf(N_\eps)) - |M_\eps| < \delta$
	Теперь наконец начнём проверять критерий для $\vf(\D)\tr M_\eps = (\vf(\D)\backslash M_\eps)\cup( M_\eps\backslash\vf(\D))$:
	\[
		\vf(\D)\backslash M_\eps = (\vf(\D)\backslash N_\eps)\cup(N_\eps\backslash M_\eps)
	\]
	А также
	\[
		M_\eps\backslash\vf\subset \vf(N_\eps)\backslash\vf
	\]
	Тогда 
	\[
		\vf(\D)\tr M_\eps \subset (\vf(\D)\tr N_\eps)\cup(N_\eps\backslash M_\eps)
	\]
	И наконец
	\[
	\mu^*(\vf(\D)\tr M_\eps) \leq \mu^*(\vf(\D)\tr N_\eps)+\mu^*(N_\eps\backslash M_\eps) < (2Mn)^n\eps + \delta
	\]
	Что завершает использование критерия.
\end{proof}

\subsection{Измеримые функции}

\begin{note}
	В этом параграфе мы будем использовать только меру Лебега, поэтому слово <<Лебега>> будем опускать.
\end{note}

\begin{definition}
	\textit{Лебеговым множеством функции} $f \colon E \to \R,\ E \subset \R^n$ называется множество
	\[
		E_a(f) = \{x \in E \colon f(x) < a\},\ a \in \R
	\]
\end{definition}

\begin{definition}
	Функция $f \colon E \to \R,\ E \subset \R^n$ и $E$ - измеримое множество, называется \textit{измеримой}, если $\forall a \in \R$ множество $E_a(f)$ измеримо.
\end{definition}

\begin{theorem} (Эквивалентное определение измеримости функций)
	Совокупность измеримых функций не изменится, если в определении лебегова множества заменить знак $<$ на любой из знаков $\le$, $\ge$, $>$.
\end{theorem}

\begin{proof}
	Знаки $\ge$ и $<$, а также $\le$, $>$ дают эквивалентные определения, так как, например, в первой паре есть связь (а дополнение измеримо тогда и только тогда, когда измеримо исходное множество)
	\[
		\{x \in E \colon f(x) \ge a\} = E \bs E_a(f)
	\]
	Докажем, что измеримость $\{x \in E \colon f(x) \le a\}$ совпадает с измеримостью $E_a(f)$. Заметим следующие равенства:
	\[
		\{x \in E \colon f(x) \le a\} = \bigcap_{m = 1}^\infty E_{a + \frac{1}{m}}(f) \quad E_a(f) = \bigcup_{m = 1}^\infty \left\{x \in E \colon f(x) \le a - \frac{1}{m}\right\}
	\]
	Случай справа (когда мы принимает измеримость с $\le$ при любом $a \in \R$) очевиден из свойств мер. Про левый же отметим одну вещь:
	\[
		E_{a + \frac{1}{m}}(f) = \left\{x \in E \colon f(x) < a + \frac{1}{m}\right\} \supset E_{a + \frac{1}{m + 1}} = \left\{x \in E \colon f(x) < a + \frac{1}{m + 1}\right\}
	\]
	Иначе говоря
	\[
		\bigcap_{m = 1}^\infty E_{a + \frac{1}{m}}(f) = \liml_{m \to \infty} E_{a + \frac{1}{m}}(f)
	\]
	а про измеримость предела теорема тоже была.
\end{proof}

\begin{note}
	Из определения следует, что любая непрерывная на измеримом множестве $E$ функция измерима. Действительно, все лебеговы множества в таком случае являются прообразами открытых в $E$ множеств. Открытость в $E$ означает, что прообраз является пересечением $E$ и какого-то открытого множества. Оба измеримы, поэтому пересечение тоже измеримо.
\end{note}

\begin{theorem} (Арифметические операции с измеримыми функциями) \label{arithmMF}
	Имеет место 2 свойства:
	\begin{enumerate}
		\item Если $f \colon \R \to \R$ непрерывна на множестве значений измеримой на $E \subset \R^n$ функции $g$, то $f \circ g$ измерима.
		
		\item Если $f, g$ измеримы на $E$, то $f \pm g$, $f \cdot g$ и (на своей области определения) $f / g$ измеримы.
	\end{enumerate}
\end{theorem}

\begin{proof}~
	\begin{enumerate}
		\item Рассмотрим лебегово множество $E_a(f \circ g)$:
		\[
			E_a(f \circ g) = \{x \in E \colon f(g(x)) < a\}
		\]
		Коль скоро $f$ непрерывна на $g(E)$, то
		\[
			f^{-1}(-\infty; a) = G \cap g(E), \text{ где $G$ - открытое множество}
		\]
		По теореме о структуре открытых множеств $G = \bscup_{m = 1}^\infty \gJ_m$ - объединение промежутков. Отсюда прообраз $f^{-1}(-\infty; a)$ можно записать так:
		\[
			f^{-1}(-\infty; a) = \bscup_{m = 1}^\infty (\gJ_m \cap g(E))
		\]
		Наконец, распишем прообраз $(f \circ g)^{-1}(-\infty; a)$:
		\[
			(f \circ g)^{-1}(-\infty; a) = \{x \in E \colon g(x) \in f^{-1}(-\infty; a)\} = \bscup_{m = 1}^\infty \{x \in E \colon g(x) \in \gJ_m\}
		\]
		Каждое множество в объединении будет измеримо, коль скоро представимо через счётное число операций над $E_a(g)$. Значит, дизъюнктное объединение тоже измеримо, что устанавливает измеримость $f \circ g$.
		
		\item Распишем лебегово множество для $f + g$:
		\begin{multline*}
			E_a(f + g) = \{x \in E \colon f(x) + g(x) < a\} = \{x \in E \colon f(x) < a - g(x)\} =
			\\
			\bigcup_{r \in \Q} \{x \in E \colon f(x) < r \wedge r < a - g(x)\} = \bigcup_{r \in \Q} E_r(f) \cap E_{a - r}(g)
		\end{multline*}
		Для умножения достаточно помнить трюк (квадрат является непрерывной функцией, а тут у нас композиция и можно применить уже доказанную часть):
		\[
			(f \cdot g)(x) = \frac{(f(x) + g(x))^2 - (f(x) - g(x))^2}{4}
		\]
		Для деления аналогично: композиция $1 / x$ - непрерывной на области определения функции и измеримой $g$.
	\end{enumerate}
\end{proof}

\begin{definition}
	\textit{Характеристической функцией} (или же \textit{индикатором}) множества $A$ называется следующая функция:
	\[
		\chi_A(x) = \System{
			&{1,\ x \in A}
			\\
			&{0,\ x \notin A}
		}
	\]
\end{definition}

\begin{proposition}
	$E \subset \R^n$ - измеримо $\lra$ $\chi_E$ - измеримая функция
\end{proposition}

\begin{proof}
	Действительно, посмотрим на то, чем является $E_a(\chi_E)$ в зависимости от $a$:
	\[
		E_a(\chi_E) = \System{
			&{\R^n,\ a > 1}
			\\
			&{\R^n \bs E,\ 0 < a \le 1}
			\\
			&{\emptyset, a \le 0}
		}
	\]
	Из этой записи уже очевидна эквивалентность условий.
\end{proof}

\begin{example} (Канторова лестница)
	Напомним вид канторова множества:
	\[
		C = [0; 1] \bs \bigcup_{m = 1}^\infty \gJ_m
	\]
	где $\gJ_1 = (1/2; 2/3), \gJ_2 = (1/9; 2/9) \sqcup (7/9; 8/9), \ldots$
	
	Определим функцию $\overline{\phi}$ в концах интервалов из $\gJ_m$, то есть на подмножестве $E = \{p/3^m \colon 0 \le p \le 3^m, m \in \N\}$.
	
	Построение будет итеративным и довольно простым. Определим $\overline{\phi}(0) = 0$ и $\overline{\phi}(1) = 1$. Далее первый интервал $\gJ_1$ сейчас между нулём и единицей, потому мы ставим там значение равное полусумме - $\frac{1}{2}$. Далее смотрим на $\gJ_2$ левый интервал в этом множестве между значениями $0$ и $\frac{1}{2}$, а правый между $\frac{1}{2}$ и $1$. Опять ставим на эти интервалы значения полусумм соседей - весь левый отображается в $\frac{1}{4}$, а правый в $\frac{3}{4}$. И так далее.
	
	Её множество значений - это $F = \{q / 2^m \colon 0 \le q \le 2^m,\ m \in \N\}$. Теперь, мы можем определить непосредственно $\phi(x)$ как функцию следующего вида:
	\[
		\phi(x) = \sup_{y \in E,\ y \le x} \overline{\phi}(y)
	\]
	$\phi(x)$ называется \textit{канторовой лестницой}, по понятным причинам. Исходя из определения, она является неубывающей. Более того, непрерывной. Это можно доказать от противного: пусть $x_0 \in [0; 1]$ - точка разрыва первого рода. Тогда множество значений $\phi$ не содержит значений из $(\phi(x_0 - 0); \phi(x_0 + 0)) \bs \{\phi(x_0)\}$. Но $F$ всюду плотно - противоречие.
	
	А что можно сказать про дифференцируемость канторовой лестницы? Мы уже знаем, что $\mu(C) = 0$. Значит, $\mu([0; 1] \bs C) = 1$. При этом
	\[
		\forall x \in [0; 1] \bs C \quad \phi'(x) = 0
	\]
	По математическим определениям это означает, что $\phi(x)$ дифференцируема \textit{почти всюду} (то есть за исключением множества нулевой меры).
\end{example}

\begin{example} (Контрпример к пункту 1 теоремы \ref{arithmMF})
	Существуют измеримая функция $f \colon \R \to \R$ и непрерывная $g \colon E \to \R$ такие, что $f \circ g$ неизмерима.
	
	Рассмотрим $\psi(x) = \frac{1}{2}(\phi(x) + x)$, где $\phi(x)$ - канторова лестница. Несложно понять, что это непрерывная, возрастающая на $[0; 1]$ биекция. Посмотрим на $\psi([0; 1] \bs C)$, в частности на меру этого образа:
	\[
		\forall x \in \ps{\frac{1}{3}; \frac{2}{3}} \Ra \psi(x) \in \ps{\frac{5}{12}; \frac{7}{12}}
	\]
	Была мера отрезка $1/3$, а стала $1/6$. Аналогично со всеми остальными имеем:
	\[
		\mu(\psi([0; 1] \bs C)) = \frac{1}{2}\mu([0; 1] \bs C) = \frac{1}{2}
	\]
	Это даёт нам возможность оценить меру $\mu(\psi(C))$:
	\[
		\mu(\psi(C)) = \mu([0; 1] \bs \psi([0; 1] \bs C)) = \frac{1}{2}
	\]
	То есть $\psi(C)$ - это измеримое множество положительной конечной меры. Значит, существует $B \subset \psi(C)$ - неизмеримое подмножество. Наша ставка заключается в том, чтобы получить $f \circ g = \chi_B$, соблюдая при этом все условия. Возьмём за $g := \psi^{-1}$, а за $f$ тогда $\chi_A$, где $A = g(B)$. По теореме об обратной функции про $g$ всё ясно, а вот измеримость $f$ следует из того факта, что $0 \le \mu(A) \le \mu(C) = 0$.
\end{example}

\begin{theorem}
	Если $\{f_m\}_{m = 1}^\infty$ - последовательность измеримых на $E \subset \R^n$ функций и при этом
	\[
		\forall x \in E\ \exists \liml_{m \to \infty} f_m(x) = f(x)
	\]
	то $f(x)$ тоже измерима на $E$.
\end{theorem}

\begin{proof}
	Покажем, что из измеримости $\forall m \in \N\ f_m$ на $E$ следует измеримость функции $\phi_t(x)$ следующего вида:
	\[
		\phi_t(x) := \sup_{m \ge t} f_m(x)
	\]
	Действительно, воспользуемся одним из эквивалентных определений измеримой функции и посмотрим на лебегово множество:
	\[
		\{x \in E \colon \sup_{m \ge t} f_m(x) \le a\} = \bigcap_{m = t}^\infty \{x \in E \colon f_m(x) \le a\} \text{ - измеримо}
	\]
	Отсюда функция, определяемая как $g(x) = \varlimsup_{m \to \infty} f_m(x) = \inf_{m \in \N} \sup_{k \ge m} f_k(x)$ тоже измерима и равна $f(x)$.
\end{proof}

\begin{definition}
	\textit{Ступенчатой функцией} будем называть измеримую функцию, имеющую конечное или счётное множество конечных значений.
\end{definition}

\begin{note}
	Иначе говоря, ступенчатая функция $h(x)$ - это такая, которая представима в следующем виде:
	\[
		h(x) = \sum_{k = 1}^\infty c_k \chi_{E_k}(x),\ E_k \text{ - измеримо}
	\]
\end{note}

\begin{theorem}
	Если $f \colon E \to [0; +\infty)$ измерима на множестве $E \subset \R^n$, то она представима пределом равномерно сходящейся на $E$ последовательности ступенчатых функций $\{h_m(x)\}_{m = 1}^\infty$ неубывающей (невозрастающей) для любого $x \in E$.
\end{theorem}